\begin{abstract}
Safety fixes for open-weight models are evaluated on the checkpoint their authors release, but users fine-tune,
compress and re-template released models, or remove refusals on purpose. We run six public safety fixes for Tiny
Aya Global, the base model and a defense of our own through one battery of six such post-release changes, with two
judges, bootstrap CIs and paired tests. Most fixes do not survive. One QLoRA epoch on \V{cfg.a2.harmful} harmful
examples takes every fix from \V{lg.A0.fixes.min}--\V{lg.A0.fixes.max}\% to \V{lg.A2.min}--\V{lg.A2.max}\% harmful
replies (Llama Guard; WildGuard \V{wg.A2.min}--\V{wg.A2.max}\%), and abliteration undoes most of them, while
capability barely moves: attacked models are unsafe, not broken. Safety at release does not predict durability
(Spearman $\rho=\V{rho.a0ret.lg}$, Llama Guard). The tools that measure safety can also mislead. Against
\V{hum.nlabels} blind hand labels, Llama Guard agrees with a human at $\kappa=\V{kappa.lg.all}$ (WildGuard
\V{kappa.wg.all}). Judging through back-translation rates \V{mlmean.A4q3.yo.wg}\% of Yoruba replies harmful after
3-bit compression (WildGuard; English \V{mlmean.A4q3.en.wg}\%), yet the author, a native Yoruba speaker, confirms
only \V{humjudge.flag.yes} of \V{humjudge.flag.n} flagged replies. Our layer-targeted defense is safer at release but
less durable. Safety in the weights is free to serve; a prompt guard cuts throughput by \V{cost.L4.RL.loss}\%. We
release the battery as an \texttt{aftermarket audit} command.
\end{abstract}
